% TOP_PROBLEM: {"id": 1859, "title": "Infinitude of Sophie Germain Primes", "category_id": 1, "category": {"id": 1, "name": "number_theory", "display_name": "Number Theory", "description": "Properties of integers, prime numbers, Diophantine equations.", "slug": "number-theory", "order_index": 1, "created_at": "2026-07-31T15:26:25.670Z"}, "status": "open"}
% FIRST_POSTED: null
% ENTRY_KIND: statement_only
% Imported from: batch08.tex; UnsolvedMath ID 1859
\documentclass[11pt]{article}
\usepackage[margin=25mm]{geometry}
\usepackage{fontspec}
\setmainfont{Latin Modern Roman}
\usepackage{amsmath,amssymb,mathtools}
\usepackage{unicode-math}
\setmathfont{Latin Modern Math}
\usepackage{xurl,hyperref}
\hypersetup{colorlinks=true,urlcolor=blue}
\setcounter{secnumdepth}{0}
\setlength{\parindent}{0pt}
\setlength{\parskip}{5pt}
\setlength{\emergencystretch}{3em}
\newcommand{\sref}[2]{\hyperlink{source:#1:#2}{[S#2]}}
\newcommand{\cataloguescope}[1]{\paragraph{Recorded target.}#1}
\begin{document}
% UnsolvedMath ID 1859; source catalogue code GUY-A7a
\setcounter{section}{1858}
\section{Infinitude of Sophie Germain Primes}\label{1859}
{\small\sffamily UnsolvedMath ID 1859 · Number Theory}
\subsection{Definitions and mathematical statement}

\paragraph{Shared notation.}$\mathbb N=\{1,2,\ldots\}$, and $\mathbb Z,\mathbb Q,\mathbb R,\mathbb C$ denote the integers, rationals, reals and complex numbers. Any other conventions are specified locally.


A prime is a positive integer greater than one whose only
positive divisors are one and itself. A Sophie Germain prime is a prime
\(p\) for which \(2p+1\) is also prime. Write
\[
 S(x)=\#\{p\le x:p\text{ and }2p+1\text{ are prime}\}.
\]
\paragraph{Statement recorded in the supplied source.}
Are there infinitely many Sophie Germain primes?
Equivalently, is \(S(x)\to\infty\) as \(x\to\infty\), or does every
bound \(B\in\mathbb N\) admit a prime \(p>B\) with \(2p+1\) prime? \sref{1859}{2}
\subsection{Short English statement}
Are there infinitely many primes whose double plus one is also prime?
\subsection{Sources}
\begin{description}
\item[\hypertarget{source:1859:1}{S1}] ULAM.AI and the UnsolvedMath Contributors, \emph{UnsolvedMath: A Curated Collection of Open Mathematics Problems}, record 1859 (GUY-A7a). CC BY 4.0. \url{https://huggingface.co/datasets/ulamai/UnsolvedMath}.
\item[\hypertarget{source:1859:2}{S2}] Matteo Bordignon and Ethan Simpson Lee, \emph{Explicit upper bounds for the number of primes simultaneously representable by any set of irreducible polynomials}, arXiv:2211.11012 (2022), Section 1, the two linear forms $x,2x+1$ and Sophie Germain primes. \url{https://arxiv.org/abs/2211.11012}.
\end{description}
\label{end:1859}
\end{document}
